\section{Constructible-hierarchy foundations}

\begin{annexitem}{annex:p06.ordinal-foundation-facts}{thm:p06-ordinal-foundation-facts}{Ordinal comparison, bounded ordinalhood, and the indexing sets}
\begin{proof}
Foundation excludes \(x\in x\): applied to the singleton \(\{x\}\),
it would otherwise have no member disjoint from that singleton.
We use this consequence only in the ambient theory \(\mathsf{ZF}\).

\paragraph{Members and comparisons.}
Let \(x\in\alpha\), where \(\alpha\) is an ordinal in the original
sense.  Transitivity gives \(x\subseteq\alpha\), so the strict
well-order on \(\alpha\) restricts to a strict well-order on \(x\).
If \(u\in v\in x\), then \(x,v,u\) all belong to \(\alpha\):
the assertions for \(v\) and \(u\) follow successively from
transitivity of \(\alpha\).  Transitivity of the strict relation
\(\in\) on \(\alpha\) gives \(u\in x\).  Thus \(x\) is transitive
and is an ordinal.

Suppose \(\gamma\subsetneq\delta\) are ordinals.  Separation on
\(\delta\) forms \(D=\delta\setminus\gamma\).  Let \(\xi\) be its
least member in the well-order of \(\delta\).  If \(\eta\in\xi\),
then \(\eta\in\delta\), and \(\eta\notin D\) by minimality of
\(\xi\).  Hence \(\eta\in\gamma\), proving
\(\xi\subseteq\gamma\).  Conversely, for \(\eta\in\gamma\), both
\(\eta,\xi\) belong to \(\delta\).  They are unequal because
\(\xi\notin\gamma\).  The alternative \(\xi\in\eta\) would put
\(\xi\) in \(\gamma\) by transitivity of \(\gamma\).
Linearity on \(\delta\) therefore gives \(\eta\in\xi\).
Thus \(\gamma\subseteq\xi\), and extensionality yields
\(\gamma=\xi\in\delta\).

For arbitrary ordinals \(\alpha,\beta\), Separation on \(\alpha\)
forms \(\gamma=\alpha\cap\beta\).  The intersection is transitive,
and its membership relation is a strict well-order inherited from
\(\alpha\); hence it is an ordinal.  If it were a proper subset
of both inputs, the preceding calculation would give
\(\gamma\in\alpha\cap\beta=\gamma\), a contradiction.  Consequently
one input includes the other, and proper inclusion is membership.
Membership implies inclusion by transitivity.  This proves the
displayed equivalence and at least one comparison alternative.
Equality together with either membership would give a self-member.
If \(\alpha\in\beta\in\alpha\), transitivity of \(\alpha\) gives
\(\alpha\in\alpha\).  Thus the alternatives are mutually exclusive.

\paragraph{The bounded predicate.}
An ordinal satisfies \(\operatorname{Tr}\), all its members are
transitive by the first paragraph, and its order is linear.  It
therefore satisfies \(\operatorname{Ord}_{\mathrm b}\).
Conversely, suppose that predicate holds of \(x\).  Membership on
\(x\) is irreflexive by Foundation and is linear by the last
conjunct.  If \(u\in v\in w\) with \(u,v,w\in x\), transitivity
of \(w\), supplied by the second conjunct, gives \(u\in w\).
Finally let \(B\subseteq x\) be nonempty.  Foundation applied to
the set \(B\) supplies \(b\in B\) with no member in \(B\).
For every distinct \(a\in B\), linearity gives \(a\in b\) or
\(b\in a\); the first is excluded.  Thus \(b\) is the least
member of \(B\).  Membership strictly well-orders \(x\), and
the first conjunct says that \(x\) is transitive.  This is exactly
the original definition of ordinal.

\paragraph{Successors, unions, and limits.}
The empty set is transitive and is vacuously strictly well-ordered.
Pairing and Union form \(\alpha+1=\alpha\cup\{\alpha\}\).
It is transitive: a member is either \(\alpha\) or a member of
\(\alpha\), and in either case all of its members lie in
\(\alpha\).  Membership extends the strict order on \(\alpha\)
by the new greatest member \(\alpha\).  For a nonempty subset of
\(\alpha+1\), its intersection with \(\alpha\), if nonempty,
supplies the least member; otherwise the subset is \(\{\alpha\}\).
This proves that the successor is an ordinal.

Let \(T\) be a set of ordinals and let \(u=\bigcup T\), a set
by Union.  If \(a\in b\in u\), some \(\delta\in T\) contains
\(b\), and its transitivity puts \(a\) in \(\delta\), hence in
\(u\).  Every member of \(u\) is an ordinal because it belongs
to a member of \(T\).  Its members are therefore transitive,
and any two are comparable by the comparison result just proved.
The bounded-predicate characterization now shows that \(u\) is
an ordinal.

Every nonzero ordinal \(\delta\) contains zero.  Indeed, its least
member \(m\) can have no member: any such member would lie in
\(\delta\) and precede \(m\).  Thus \(m=\varnothing\).
For a limit ordinal \(\lambda\), transitivity gives
\(\bigcup\lambda\subseteq\lambda\).  If this inclusion were
proper, \(\mu=\bigcup\lambda\) would be an ordinal member of
\(\lambda\).  Every \(\beta\in\lambda\) is a subset of \(\mu\),
so the comparison equivalence gives \(\beta\leq\mu\).  It follows
that \(\lambda\subseteq\mu+1\).  Conversely, \(\mu\in\lambda\)
and transitivity give \(\mu+1\subseteq\lambda\).
Then \(\lambda=\mu+1\), contrary to the definition of limit.
Hence \(\bigcup\lambda=\lambda\).
If \(\beta<\lambda\), then \(\beta+1\subseteq\lambda\), so
\(\beta+1\leq\lambda\).  Equality would make \(\lambda\) a
successor.  Therefore \(\beta+1<\lambda\).

\paragraph{The natural-number index set.}
The earlier Separation construction makes \(\omega\) the least
inductive set.  Separate from \(\omega\) the set of all \(n\)
such that \(n\) is an ordinal, \(n\subseteq\omega\), and \(n\)
is either zero or a successor ordinal.  This subset contains zero.
If it contains \(n\), it contains \(n+1\): the successor is an
ordinal, belongs to \(\omega\) by inductivity, and is included
in \(\omega\) because \(n\subseteq\omega\) and \(n\in\omega\).
The subset is inductive, and leastness makes it all of \(\omega\).
Thus \(\omega\) is transitive, consists of transitive ordinals,
and has linear membership by comparison of its members.
The bounded-predicate characterization proves it is an ordinal.
The same argument proved that its nonzero members are successors.
Induction on \(\omega\) follows either from its least inductivity
by Separation or from transfinite induction restricted to it.
\end{proof}
\end{annexitem}

\begin{annexitem}{annex:p06.definable-ordinal-recursion}{thm:p06-definable-ordinal-recursion}{The approximation construction and its exact set bounds}
\begin{proof}
All occurrences of \(G\) in this proof are replaced by the fixed
formula \(\vartheta\) of
\cref{def:p06-recursion-approximation}.  Consequently the properties
used in Separation, Replacement, and transfinite induction are
formulas of \(\Lang_{\in}\) with the fixed set parameters \(\bar p\).

For a set function \(f\) and an ordinal \(\xi\), its restriction
is formed by Separation on \(f\), retaining the pairs whose first
coordinate belongs to \(\xi\).  If \(\xi\subseteq\dom(f)\), the
result has domain \(\xi\).  Taking a second restriction gives
\((f\mathbin{\upharpoonright}\xi)\mathbin{\upharpoonright}\eta=f\mathbin{\upharpoonright}\eta\) for \(\eta\leq\xi\):
both graphs retain exactly the pairs with first coordinate in
\(\eta\).  It follows directly from the defining universal
condition that a restriction of an approximation is an approximation.

\paragraph{Compatibility before existence.}
Let \(f,g\) be any approximations, with domains \(\gamma,\delta\).
Let \(\tau\) be the smaller of these ordinals.  Separation on
the set \(\tau\) forms
\[
 D=\{\xi\in\tau:f(\xi)\ne g(\xi)\}.
\]
If \(D\ne\varnothing\), let \(\mu\) be its least member.
The restrictions \(f\mathbin{\upharpoonright}\mu,g\mathbin{\upharpoonright}\mu\) have the same domain
and agree at every input, so their graphs are equal.  The
approximation clauses and functionality of \(\vartheta\) imply
\(f(\mu)=g(\mu)\), a contradiction.  Thus the approximations
agree on their common domain.  In particular there is at most one
approximation of any specified length.

\paragraph{Existence at each ordinal.}
Apply the completed transfinite induction theorem
\cref{thm:p06-transfinite-induction} to
\(\exists f\,\operatorname{App}_G(\gamma,f)\).
At zero the empty function works.  At \(\gamma=\alpha+1\), let
\(f_\alpha\) be the earlier approximation, and let \(y\) be the
unique value specified by \(\vartheta(f_\alpha,y,\bar p)\).
Pairing and Union form the graph
\[
 f_\alpha\cup\{\langle\alpha,y\rangle\}.
\]
The new first coordinate is outside the old domain, so this graph
is a function on \(\alpha+1\).  Its restriction to \(\alpha\)
is \(f_\alpha\).  Old recursion equations remain true, and the
new one is precisely the choice of \(y\).

Let \(\lambda\) be nonzero and limit.  For each \(\beta\in\lambda\),
the induction hypothesis and the compatibility argument give
exactly one \(g\) satisfying \(\operatorname{App}_G(\beta,g)\).
The Replacement instance with domain \(\lambda\) and defining
formula \(\operatorname{App}_G(\beta,g)\) therefore forms the set
\[
 H=\{g:\exists\beta\in\lambda\;
                    \operatorname{App}_G(\beta,g)\}.
\]
Union supplies \(f=\bigcup H\).  Compatibility ensures that
two pairs in \(f\) with the same first coordinate have the same
second coordinate.  Its domain is
\(\bigcup_{\beta<\lambda}\beta=\lambda\), using
\cref{thm:p06-ordinal-foundation-facts}.
For \(\xi\in\lambda\), the ordinal \(\xi+1\) also belongs to
\(\lambda\).  The member \(f_{\xi+1}\) of \(H\) contains its
value at \(\xi\).  This graph is included in the union \(f\),
which is already a function, so
\(f\mathbin{\upharpoonright}(\xi+1)=f_{\xi+1}\); hence its recursion equation at
\(\xi\) is the equation for \(f\).  This proves existence at
\(\lambda\) and completes the induction.  No collection over
all ordinals was formed: the Replacement domain was the set
\(\lambda\) at that stage.

The compatibility result now gives
\(f_\beta=f_\gamma\mathbin{\upharpoonright}\beta\) whenever \(\beta\leq\gamma\).
It also shows that the displayed definition of \(F\) has exactly
one value at each ordinal and that its restriction to \(\gamma\)
has graph \(f_\gamma\).  Its defining equation follows from the
approximation of length \(\alpha+1\).
For any set \(D\) of ordinals, Replacement on \(D\), with the
functional formula assigning
\(\langle\alpha,y\rangle\) where \(F(\alpha)=y\), produces the
graph of \(F\mathbin{\upharpoonright} D\).
Finally, any competing definable class solution has set
restrictions by this same Replacement argument.  Its restriction
to \(\alpha+1\) is an approximation and therefore equals
\(f_{\alpha+1}\).  The two solutions agree at \(\alpha\), for
every ordinal \(\alpha\).  This is uniqueness of the definable
class function, expressed as a scheme for its defining formulas.
\end{proof}
\end{annexitem}

\begin{annexitem}{annex:p06.set-sized-definability}{thm:p06-set-sized-definability}{A set of assignments and finite truth computations}
\begin{proof}
\paragraph{The bounding sets.}
For sets \(a,b\), a Kuratowski pair \(\langle x,y\rangle\) with
\(x\in a,y\in b\) belongs to \(\mathcal{P}(\mathcal{P}(a\cup b))\): each
of its members is a subset of \(a\cup b\).  Separation on this
double power set, with the condition
\(\exists x\in a\,\exists y\in b\;(z=\langle x,y\rangle)\),
constructs \(a\times b\).  Repeating this construction gives the
set \(B=\omega\times(\omega\times\omega)\) of triple records,
where \((r,i,j)\) denotes \(\langle r,\langle i,j\rangle\rangle\).
Every finite list of records is a member of
\(\mathcal{P}(\omega\times B)\).  Separation on this power set, using
the function condition and
\(\exists n\in\omega\;(\dom(c)=n\land n\ne0)\), forms all
nonempty finite lists.  A further Separation applies the nine
explicit record conditions in
\cref{def:p06-finite-formula-instructions}, with earlier positions
bounded by the current position.  Its result is \(\mathcal C\).
This defines a set by a single finite formula; no collection of
all set-theoretic formulas is being assumed as a proper class.

For a fixed nonempty \(A\), Separation on \(\mathcal{P}(\omega\times A)\)
with the function and domain conditions gives \(S_A\).
Choose one \(a_0\in A\).  The graph
\(\{\langle i,a_0\rangle:i\in\omega\}\), formed by Separation
on \(\omega\times A\), is a member of \(S_A\).
This proves nonemptiness without a choice function on a family of
sets.  An update \(s[i\mapsto a]\) is the union of the singleton
\(\{\langle i,a\rangle\}\) with the subset of \(s\) whose first
coordinate is different from \(i\).  Separation, Pairing, and
Union give this set, and its domain and range show that it lies
in \(S_A\).

\paragraph{Truth sets for one code.}
Fix \(c\in\mathcal C\) of length \(n\), a nonzero member of
\(\omega\).  At each position \(k<n\), construct a subset
\(T_k\) of \(S_A\).  With the record at \(k\) as in the body,
the nine clauses are
\[
\begin{array}{c|l}
0&T_k=\{s\in S_A:s(i)=s(j)\}\\
1&T_k=\{s\in S_A:s(i)\in s(j)\}\\
2&T_k=\varnothing\\
3&T_k=\{s\in S_A:s\notin T_i\}\\
4&T_k=\{s\in S_A:s\in T_i\land s\in T_j\}\\
5&T_k=\{s\in S_A:s\in T_i\lor s\in T_j\}\\
6&T_k=\{s\in S_A:s\notin T_i\lor s\in T_j\}\\
7&T_k=\{s\in S_A:\forall a\in A\;
                         s[i\mapsto a]\in T_j\}\\
8&T_k=\{s\in S_A:\exists a\in A\;
                         s[i\mapsto a]\in T_j\}.
\end{array}
\]
Every right side is a Separation instance on the already available
set \(S_A\).  To justify the simultaneous finite construction
without assuming another recursion theorem, define an operation
on every ordinal-domain function \(h\) as follows.  If
\(\delta=\dom(h)\) belongs to \(n\), use the applicable clause
for \(c(\delta)\), replacing earlier \(T_i,T_j\) by
\(h(i),h(j)\).  Every referenced earlier index belongs to
\(\delta\).  If \(\delta\notin n\), return \(\varnothing\).
This is a total formula-defined operation on such functions,
even when their values are arbitrary sets.  Each branch returns
a unique subset of \(S_A\).
Apply \cref{thm:p06-definable-ordinal-recursion} and restrict its
solution to \(n\).  The resulting set function
\(t=\langle T_k:k<n\rangle\) exists and satisfies all clauses.
Finite induction proves uniqueness, since the clause at \(k\)
uses only earlier values.

The assertion that \(t\) is this computation is a single
first-order formula in \(A,c,t\).  It says that \(t\) has domain
\(n\), that every \(t(k)\) is a subset of \(S_A\), and, for
each \(k\in n\), that one of the nine record-and-membership
clauses above holds for all \(s\in S_A\).  The symbols \(S_A\),
updates, and record projections have the set definitions already
given.  Since \(n\ne0\), it is a successor in \(\omega\), so
its final position is the unique \(m\) with \(n=m+1\).
For unrestricted set arguments, define
\[
\begin{aligned}
 \operatorname{Sat}(A,c,s)
 \quad:\Longleftrightarrow\quad
 &A\ne\varnothing\land c\in\mathcal C\land s\in S_A\\
 &{}\land\exists n\in\omega\,\exists m\in n\,\exists t\;
 \bigl(\dom(c)=n\land n=m+1\\
 &\hspace{13mm}\land
 t\text{ is the displayed computation for }A,c
                  \land s\in t(m)\bigr).
\end{aligned}
\]
Separation on \(\mathcal C\times S_A\) forms
\(\operatorname{Sat}_A\).  Existence and uniqueness of each
computation prove its existence and uniqueness with the required
clauses.  The displayed existential assertion, with its finite
defining abbreviations expanded, is the promised uniform formula.
Its parameter \(A\) is a set; the construction has not supplied
a satisfaction predicate for the universe.

\paragraph{Agreement with formulas and with \(\operatorname{Def}\).}
For every code, induction over its finite positions unfolds each
record to a membership formula.  The atomic, falsehood, Boolean,
and quantified rows are exactly the clauses of
\cref{sec:fol-semantics-02}.  The induction therefore identifies
\(T_k\) with the satisfaction set of the unfolded formula at
position \(k\).  Reusing a position copies that same formula
and truth set into both occurrences.
Conversely, process a formula from its subformulas to its root.
Atoms and falsehood require one record.  For a unary constructor,
append its record to the child's list.  For a binary constructor,
concatenate the two child lists, add the length of the first list
to every position-reference field in the second list, and append
the binary record pointing to their final positions.  Variable
indices are left unchanged; only fields naming earlier positions
are shifted.  This finite shift requires no ordinal-arithmetic law:
if the first length is \(m\), define \(r(0)=m\) and
\(r(k+1)=r(k)+1\) by the proved recursion theorem.  Induction on
\(k\) shows that its successive positions are distinct, lie
outside the first list, and preserve every earlier-position
comparison.  Replacing a second-list index \(k\) by \(r(k)\)
therefore gives the claimed concatenation; Replacement on the
finite domains forms its graph.  Quantifiers append a record
containing the bound variable index and the child's final position.
Induction on the
formula proves that every referenced position is earlier and that
the output has its original satisfaction condition.  Derived
connectives and bounded quantifiers are first expanded into the
fixed logical language.

For nonempty \(A\), Separation on \(\mathcal{P}(A)\), using the formula
in the theorem statement, now gives a unique set of subsets.
Any witnessing code unfolds to a finite formula.  Only finitely
many of its variables other than \(v_0\) are free, so only the
corresponding finite tuple of values of \(s\) is relevant, by
the formula-coincidence lemma \cref{lem:fol-coincidence}.
Its subset is therefore definable with finitely many parameters
in the original sense.
Conversely, take an original defining formula and its finite
parameter tuple.  If its distinguished variable is not \(v_0\),
interchange its variable index with zero throughout the formula,
including bound occurrences, and move the parameter assignments
accordingly.  This bijective renaming preserves the truth clauses:
atomic indices are renamed together, Boolean clauses commute with
renaming, and updating a renamed bound index gives exactly the
renamed updated assignment in each quantifier clause.
Thus the distinguished variable may be taken to be \(v_0\).
Extend the finite assignments to \(s\in S_A\) by assigning
the fixed \(a_0\) to all other indices.  Separation on
\(\omega\times A\) constructs its graph.  A code for the formula
then witnesses precisely the required subset in the display.
The two definitions of \(\operatorname{Def}(A)\) agree.

For \(A=\varnothing\), Pairing forms \(\{\varnothing\}\).
Combining this branch with the nonempty branch gives one
formula defining a unique set value for every \(A\).
The empty branch does not call the satisfaction construction.
\end{proof}
\end{annexitem}

\begin{annexitem}{annex:p06.bounded-transitive-absoluteness}{thm:p06-bounded-transitive-absoluteness}{Absoluteness by bounded-formula induction}
\begin{proof}
Fix a nonempty transitive set \(A\).  For each bounded formula,
prove the assertion by induction over its finite formation.
Membership and equality of assigned members of \(A\) have the
same meaning in the structure and in the ambient formula.
Falsehood is false in both.  The induction passes through
negation, conjunction, disjunction, and implication because
their truth clauses are the same Boolean operations on the
already equivalent assertions.

For \(\exists u\in v\;\psi\), let \(b\in A\) be the value of
\(v\).  The quantified variables \(u,v\) are distinct, so this
value remains fixed when \(u\) is updated.  The structural truth
clause asks for an \(a\in A\) such that \(a\in b\) and \(\psi\)
holds under that update.  Transitivity gives \(b\subseteq A\).
The possible witnesses are therefore exactly the ambient members
of \(b\).  The induction hypothesis applies at each such update,
and the existential assertions agree.  For
\(\forall u\in v\;\psi\), the structural clause imposes \(\psi\)
on those \(a\in A\) that belong to \(b\), and no condition on
the other members of \(A\).  Again these relevant values are
exactly the members of \(b\), and the induction hypothesis gives
the ambient bounded universal assertion.
This completes the induction.  The special case for
\(\operatorname{Ord}_{\mathrm b}(x)\) follows from its bounded
syntax and its ambient characterization in
\cref{thm:p06-ordinal-foundation-facts}.  No axiom of \(\mathsf{ZF}\)
has been assumed to hold inside \(A\).
\end{proof}
\end{annexitem}

\begin{annexitem}{annex:p06.def-transitive-growth}{thm:p06-def-transitive-growth}{Inclusion and transitivity at a definability step}
\begin{proof}
Every member of \(\operatorname{Def}(A)\) is a subset of \(A\)
by its definition.  This includes the empty case, because the
sole member of \(\operatorname{Def}(\varnothing)\) is the empty
subset.  For nonempty \(A\), the formula \(v_0=v_0\) defines
all of \(A\), so \(A\in\operatorname{Def}(A)\).  For empty
\(A\), the same membership is the separate base clause.

Suppose \(A\) is transitive and \(x\in A\).  Then
\(x\subseteq A\), and the formula \(v_0\in v_1\), with
parameter \(x\) assigned to \(v_1\), defines exactly \(x\)
over \(A\).  Thus \(x\in\operatorname{Def}(A)\), proving
\(A\subseteq\operatorname{Def}(A)\); if \(A\) is empty this
inclusion is vacuous.  Finally, if
\(u\in v\in\operatorname{Def}(A)\), then \(v\subseteq A\),
so \(u\in A\subseteq\operatorname{Def}(A)\).
This is transitivity of \(\operatorname{Def}(A)\).
\end{proof}
\end{annexitem}

\begin{annexitem}{annex:p06.hierarchy-recursion-existence}{thm:p06-hierarchy-recursion-existence}{The formula defining each level}
\begin{proof}
For a function \(f\) with ordinal domain \(\delta\), define
\[
 G_L(f)=
 \begin{cases}
  \varnothing,&\delta=0,\\
  \operatorname{Def}(f(\beta)),&\delta=\beta+1,\\
  \bigcup\operatorname{ran}(f),&\delta\text{ is nonzero and limit}.
 \end{cases}
\]
These cases are exhaustive and disjoint.  The predecessor in the
successor case is unique: comparison of two proposed predecessors
and transitivity would otherwise make the smaller successor a
proper subset of the larger one.  More explicitly, if
\(\beta<\gamma\), then \(\beta+1\subseteq\gamma\), whereas
\(\gamma\in\gamma+1\) and \(\gamma\notin\gamma\); hence
\(\beta+1\ne\gamma+1\).
The range in the limit case is a set by Replacement on
\(\delta\) with the functional formula assigning \(f(\xi)\)
to \(\xi\); Union gives its union.
The nonempty and empty definitions of \(\operatorname{Def}\)
are uniformly functional by \cref{thm:p06-set-sized-definability}.
Thus this display expands to a single formula defining one set
for every ordinal-domain function.

Apply \cref{thm:p06-definable-ordinal-recursion} to \(G_L\).
Its zero value is empty.  At \(\beta+1\), its value is
\(\operatorname{Def}(L_\beta)\).  At a nonzero limit, its value
is the union of the earlier values.  These are exactly the
previously stated clauses for \(L\), so uniqueness identifies
the hierarchy with that definition.
Replacing \(\operatorname{Def}\) by \(\mathcal{P}\) gives another
total definable operation \(G_V\), since Power Set supplies its
unique successor value.  The same recursion theorem supplies
exactly the previously stated \(V\)-hierarchy clauses.
This proves their existence and the assertion about set
restrictions without using the earlier sketched hierarchy theorem.

For clarity, if \(\Lambda(\alpha,z)\) is the formula from the
recursion theorem saying that the \(G_L\) value at \(\alpha\)
is \(z\), then the expanded constructibility formula is
\[
 \exists\alpha\,\exists z\;
    (\operatorname{Ord}_{\mathrm b}(\alpha)
                   \land\Lambda(\alpha,z)\land x\in z).
\]
It quantifies over sets.  Its reference to \(\operatorname{Def}\)
inside the approximation formula uses the set-sized construction
above, so this definition introduces no universe truth predicate.
\end{proof}
\end{annexitem}

\begin{annexitem}{annex:p06.constructible-level-height}{thm:p06-constructible-level-height}{Transitivity, increasing levels, and their exact ordinal members}
\begin{proof}
All induction assertions below are first-order formulas after
expanding the definable hierarchy of
\cref{thm:p06-hierarchy-recursion-existence}.  Thus the completed
transfinite induction theorem applies to each of them.

\paragraph{Transitivity and inclusion.}
The empty level is transitive.  At a successor, transitivity of
the previous level and \cref{thm:p06-def-transitive-growth}
give transitivity of its definability set.  At a nonzero limit
\(\lambda\), if \(u\in v\in L_\lambda\), there is a
\(\beta<\lambda\) with \(v\in L_\beta\).  Transitivity of
that earlier level puts \(u\) in \(L_\beta\), hence in
\(L_\lambda\).  This proves transitivity at all levels, and
the growth lemma consequently gives
\(L_\alpha\subseteq L_{\alpha+1}\) for all \(\alpha\).

Induct now on \(\beta\) to show that
\(L_\alpha\subseteq L_\beta\) for every \(\alpha\leq\beta\).
Equality of indices needs no argument.  If \(\beta=\gamma+1\)
and \(\alpha<\beta\), then \(\alpha\leq\gamma\), so the
induction hypothesis followed by the successor inclusion gives
the assertion.  If \(\beta\) is limit and \(\alpha<\beta\),
the set \(L_\alpha\) is one of the sets whose union defines
\(L_\beta\), giving the inclusion directly.  The zero case is
the equality case.

\paragraph{Exact height, including the definability step.}
Separate the ordinal members from each set \(L_\alpha\), using
\(\operatorname{Ord}_{\mathrm b}\); the expression
\(L_\alpha\cap\operatorname{Ord}\) names this set.  At zero it
is empty.  Suppose inductively that
\(L_\beta\cap\operatorname{Ord}=\beta\), and put \(A=L_\beta\).
If \(\beta=0\), then \(A=\varnothing\) and
\(\operatorname{Def}(A)=\{\varnothing\}\).  Its ordinal members
are exactly \(1=\{0\}\), and \(\beta=0\) belongs to it.

If \(\beta>0\), then \(0\in\beta\), so the induction
identity implies \(0\in A\).  Thus \(A\) is nonempty, and
its already established transitivity permits
\cref{thm:p06-bounded-transitive-absoluteness}.  The subset
defined over \(A\) by the parameter-free formula
\(\operatorname{Ord}_{\mathrm b}(v_0)\) is exactly
\[
 \{x\in A:(A,\in)\entails\operatorname{Ord}_{\mathrm b}(x)\}
   =\{x\in A:x\text{ is an ordinal}\}=\beta.
\]
It follows from the actual definition of \(\operatorname{Def}\)
that \(\beta\in\operatorname{Def}(A)=L_{\beta+1}\).
The preceding empty case supplies the same conclusion at zero.
In both cases every \(\xi<\beta\) already belongs to \(A\)
and therefore belongs to \(L_{\beta+1}\) by inclusion.
All ordinals at most \(\beta\) consequently occur there.

For the reverse inclusion, let an ordinal \(\eta\) belong to
\(L_{\beta+1}=\operatorname{Def}(A)\).
Then \(\eta\subseteq A\).  Each member \(\xi\) of \(\eta\)
is an ordinal and belongs to \(A\), so the induction identity
puts \(\xi\) in \(\beta\).  Hence \(\eta\subseteq\beta\),
and ordinal comparison gives \(\eta\leq\beta\).
We have proved
\(L_{\beta+1}\cap\operatorname{Ord}=\beta+1\).

For a nonzero limit \(\lambda\) and an ordinal \(\eta\),
the union clause and the induction hypotheses give
\[
\begin{aligned}
 \eta\in L_\lambda
 &\quad\Longleftrightarrow\quad
       \exists\beta<\lambda\;(\eta\in L_\beta)\\
 &\quad\Longleftrightarrow\quad
       \exists\beta<\lambda\;(\eta\in\beta)\\
 &\quad\Longleftrightarrow\quad \eta\in\lambda.
\end{aligned}
\]
The last equivalence is \(\bigcup\lambda=\lambda\), already
proved for limit ordinals.  This completes the exact-height
induction and also proved \(\alpha\in L_{\alpha+1}\) at every
index.  The ordinal \(\alpha\) cannot belong to \(L_\alpha\):
the height identity would imply \(\alpha\in\alpha\).
Thus the previously proved successor inclusion is strict.

\paragraph{Comparison with the cumulative hierarchy.}
This final induction uses only the existence and defining clauses
of \(V\) proved in \cref{thm:p06-hierarchy-recursion-existence}.
At zero both levels are empty.  If \(L_\alpha\subseteq V_\alpha\)
and \(x\in L_{\alpha+1}\), then
\(x\subseteq L_\alpha\subseteq V_\alpha\), so
\(x\in\mathcal{P}(V_\alpha)=V_{\alpha+1}\).
At a nonzero limit, any \(x\in L_\lambda\) belongs to some
\(L_\beta\) with \(\beta<\lambda\).  The induction hypothesis
puts it in \(V_\beta\), which is included in the union
\(V_\lambda\).  Hence \(L_\alpha\subseteq V_\alpha\) for all
\(\alpha\).  Neither existence of ranks nor the assertion that
every set occurs in a \(V\)-level was used.
\end{proof}
\end{annexitem}

\begin{annexitem}{annex:p06.constructible-class-foundation}{thm:p06-constructible-class-foundation}{The class conclusion from the level theorem}
\begin{proof}
If \(x\in y\in L\), there is an ordinal \(\alpha\) such that
\(y\in L_\alpha\).  Transitivity of that level gives
\(x\in L_\alpha\), hence \(x\in L\).  For every ordinal
\(\alpha\), the level theorem gives
\(\alpha\in L_{\alpha+1}\), so \(\alpha\in L\).

If \(L\) were a set \(l\), Separation on \(l\) with
\(\operatorname{Ord}_{\mathrm b}\) would form a set \(O\)
containing exactly all ordinals.  Its union \(\nu=\bigcup O\)
would be an ordinal by \cref{thm:p06-ordinal-foundation-facts}.
The ordinal \(\nu+1\) would belong to \(O\), so
\(\nu\in\nu+1\subseteq\bigcup O=\nu\), contrary to Foundation.
Thus the constructible class is proper.
\end{proof}
\end{annexitem}
